\subsection{E(b). Two commuting operator groups}

We retain the names $\theta$ and $\omega$ for the standard order-two and order-three lifts on $\widehat F_2$; Grothendieck's $\sigma_0$ and $\rho$ denote the corresponding choices on scans 316--317, \S25 bis, folios 159--160. A change of base path conjugates a lift by $\widehat F_2$. We keep this ambiguity in the statements below.

\begin{claim}{E.20}{Proved}{The two lifting conditions cut out exactly the multiplier-$1$ subgroup modulo $3$}
In $B=\Out_{\rm lac}(\widehat\pi_{0,3})$ let $S_3$ be the geometric positive subgroup. Its centralizer is $GT_0$, with the normalized description $(\lambda,f)$ of Harbater--Schneps. Every element of $GT_0$ has an automorphism representative commuting with the chosen order-two lift $\theta$. The elements with an automorphism representative commuting with the chosen order-three lift $\omega$ are exactly
\[
 GT_0^{(3)}=\{(\lambda,f)\in GT_0:\lambda\equiv1\pmod3\}.
\]
This subgroup has index two. The two representatives need not be the same representative; this is the two-condition formulation on scans 316--317.
\end{claim}
\begin{proof}
Harbater--Schneps, Main Theorem (a) and Theorem A(i) \cite{R-HS}, identify the centralizer and state precisely
\[
 [\theta,F]=\Inn(f),\qquad
 [\omega,F]=\Inn(y^mf),\qquad m=(\lambda-1)/2,
\]
for $F(x)=x^\lambda$, $F(y)=f^{-1}y^\lambda f$. Here $[a,b]=aba^{-1}b^{-1}$. A representative $\Inn(h)F$ commutes with $\theta$ precisely when
\[
 \theta(h)f=h,
\]
since $F_2$ has trivial centre. It commutes with $\omega$ precisely when $\omega(h)y^mf=h$. These are the two exact coboundary equations. We determine their solvability without imposing the pentagon equation.

The positive lift extension of the $S_3$ action on the three-punctured sphere is
\[
 E=C_2*C_3=\langle s,r\mid s^2=r^3=1\rangle,
 \qquad E\longrightarrow S_3,
\]
whose kernel is the peripheral free group $F_2$. Its cusp generator can be chosen as $e=sr$; $e^2$ is a primitive cusp in that kernel. Complete this exact sequence by D.1. An outer centralizer element, represented by $F$, acts on $\widehat E$ by conjugating its image in $\Aut(\widehat F_2)$ and induces the identity on $S_3$. It sends $e^2$ to a $\widehat F_2$-conjugate of $e^{2\lambda}$. After that conjugation it sends $e$ into $C_{\widehat E}(e^{2\lambda})=\overline{\langle e\rangle}$. To prove this equality, put $L=\overline{\langle e^2\rangle}$. A centralizing element preserves its $\widehat F_2$-conjugacy class, so its image in $S_3$ stabilizes the associated puncture. That stabilizer is $\langle q(e)\rangle$ of order two. Multiplication by $e$ when necessary puts the element in $\widehat F_2$, where Part A gives $C_{\widehat F_2}(e^{2\lambda})=L$. This proves the equality and identifies the whole centralizer with $\widehat\Z$. Doubling is injective in that group, so the image of $e$ is $e^\lambda$.

Consequently this automorphism acts on $\widehat E_{\rm ab}=C_6$ by multiplication by $\lambda$. Every order-two element is conjugate to $s$, and every order-three subgroup is conjugate to $\langle r\rangle$, by D.2 or the free profinite product theorem. A lift of the given transposition is $F_2$-conjugate to $s$: a total-group conjugator projects to $C_{S_3}(s)=\langle s\rangle$, and multiplication by $s$ adjusts it into the kernel. Thus there is one kernel-conjugacy class for the order-two lift. It is fixed by every centralizer element, proving solvability of the first equation.

There are exactly two kernel-conjugacy classes of lifts of the specified three-cycle. They are represented by $r$ and $s r^{-1}s^{-1}$. Indeed a total-group conjugator has either an even or an odd image in $S_3$; adjusting by $\langle r\rangle$ or by $s\langle r\rangle$ gives these two representatives. They are distinct because their images in the order-three quotient of $E_{\rm ab}$ are inverse nonidentity elements. The action of $(\lambda,f)$ preserves the class of $r$ exactly when $\lambda\equiv1\pmod3$, by the computed action on $C_6$. When it does, the same adjustment by $r$ puts its conjugator in the kernel. This is precisely solvability of the second coboundary equation. The argument treats elements, so it retains the label of the three-cycle.

For comparison, Lochak--Schneps, Proposition 3 and Theorem 2 \cite{R-LS}, compute the same two nonabelian cohomology sets: one class for $\theta$, two for $\omega$, with the second class selected by $\lambda\equiv-1\pmod3$. The first two parts of their proof use only the two-cycle and three-cycle equations, so their computation also applies to $GT_0$. Finally $(-1,1)\in GT_0$ exchanges the two classes. The homomorphism to $(\Z/3)^\times$ is therefore surjective and the asserted index is exactly two. The residue calculation is independently recorded in the executed program.
\end{proof}

\begin{claim}{E.21}{Proved}{The arithmetic subgroup and the two fixed points}
The intersection of the usual Galois image $G_{\Q}\to GT_0$ with $GT_0^{(3)}$ is the image of $G_{\Q(\sqrt{-3})}$. The order-three fixed points in $U=\mathbf P^1-\{0,1,\infty\}$ are $e^{\pm i\pi/3}$, and the order-two lift has one fixed point in $U$. Thus the requirement to fix the chosen order-three section is a proper finite-index restriction, even on the Galois subgroup.
\end{claim}
\begin{proof}
The multiplier of the Galois action is the cyclotomic character, by its action on the tame inertia at $0$. Its reduction modulo $3$ is the action on $\mu_3$ and has fixed field $\Q(\zeta_3)=\Q(\sqrt{-3})$. E.20 identifies precisely its kernel. The three-cycle $z\mapsto1/(1-z)$ has fixed-point equation $z^2-z+1=0$, and the transposition $z\mapsto1-z$ has fixed-point equation $z=1/2$. These points lie in $U$. Complex conjugation exchanges the former pair and fixes the latter point, which also proves properness directly.

On scan 312 the relevant subgroup is obtained by fixing a finite set of critical data; it is therefore of finite index. An assertion that every element of $G_\Q$ lifts commuting with a specified one of the order-three lifts would contradict the displayed computation. We record the scan's wording as editorial evidence, not as a settled reading.
\end{proof}

\begin{claim}{E.22}{Proved}{Normalizer, stabilizer, and the forced lift}
Let $1\to\pi\to E\xrightarrow q\mathcal G\to1$ have centre-free kernel, and let $K\leq E$ be a finite subgroup on which $q$ is injective. Assume $C_\pi(K)=1$, as holds for a nontrivial rotation stabilizer in the anabelian geometric situation and its completion. Put $G_Q=q(K)$ and let $Q$ denote the $\pi$-conjugacy class of $K$. The restriction
\[
 q:N_E(K)\longrightarrow\mathcal G
\]
is injective and has image the stabilizer $\mathcal G_Q$ of $Q$. Its inverse on that image is a canonical homomorphism. In a two-group situation $\mathcal G=\Gamma G$ with $\Gamma$ centralizing $G$, its restriction gives the forced lift of $\Gamma_Q=\Gamma\cap\mathcal G_Q$, and this lift centralizes $K$.
\end{claim}
\begin{proof}
An element of $N_E(K)\cap\pi$ acts trivially on $q(K)$, and the restriction of $q$ to $K$ is injective. It therefore centralizes $K$, so the kernel is $C_\pi(K)=1$. A quotient element stabilizes $Q$ precisely when one, and hence every, lift carries $K$ to a $\pi$-conjugate. Multiplying that lift by the appropriate element of $\pi$ produces an element of $N_E(K)$ above it. This proves the image statement and uniqueness of the lift. The inverse is a homomorphism because it is the inverse of an injective homomorphism onto its image. For $\gamma\in\Gamma_Q$, conjugation of its lift on $K$ projects to conjugation by $\gamma$ on $G_Q$, which is trivial. Injectivity on $K$ then makes this conjugation trivial. In the profinite case all these maps are continuous, and the inverse is continuous by compactness of the closed normalizer and the Hausdorff property of the quotient.

For the centralizer, the precise image is $\mathcal G_Q\cap C_{\mathcal G}(G_Q)$, rather than all of $\mathcal G_Q$. When $G_Q$ is nonabelian dihedral, its own elements already distinguish these images. For a reflection in an oval the kernel is $\Z$ or $\widehat\Z$ by Part B and D.4, so the stated injectivity is replaced by an exact sequence with that kernel. These computations identify exactly the rigidity used in the manuscript's argument.
\end{proof}

\begin{claim}{E.23}{Proved}{Explicit discrete centralizers and pointwise kernels}
For the three actions in the question, the full extended geometric groups, their centralizers, their normalizers, and their pointwise kernels on the exceptional points are as follows. Here $D_{2n}$ has order $2n$.
\[
\begin{array}{c|c|c|c|c}
 (g,\nu),G&M&C_M(G)&N_M(G)&C_M(G)^{\natural}\\ \hline
 (0,3),S_3&S_3\times C_2&C_2&S_3\times C_2&1\\
 (1,1),C_4&\mathrm{GL}_2(\Z)&C_4&D_8&C_2\\
 (1,1),C_6&\mathrm{GL}_2(\Z)&C_6&D_{12}&1
\end{array}
\]
The last column fixes every exceptional point individually, not merely its orbit under $G$.
\end{claim}
\begin{proof}
For type $(0,3)$, the geometric group is $S_3$ times complex conjugation. The exceptional points are the three transposition fixed points $-1,1/2,2$ and the two order-three fixed points of E.21. Complex conjugation fixes the first three and interchanges the last two. The centralizer of $S_3$ is therefore $C_2$ and its pointwise kernel is trivial; its normalizer is the whole product.

Use the matrices
\[
 R_4=\begin{pmatrix}0&-1\\1&0\end{pmatrix},\qquad
 R_6=\begin{pmatrix}0&-1\\1&1\end{pmatrix}
\]
on $\R^2/\Z^2$, punctured at $0$. A rational matrix commuting with either one is $aI+bR_n$. For an integral commuting matrix $a,b\in\Z$. Its determinant is $a^2+b^2$ for $R_4$ and $a^2+ab+b^2$ for $R_6$. Equating this nonnegative norm to $1$ gives four and six units respectively, proving the centralizer assertions. Conjugation in the normalizer embeds its quotient by the centralizer in $\Aut(C_n)=C_2$ for $n=4,6$. Reflection in the square or hexagonal lattice realizes inversion, proving the two normalizer assertions.

For $C_4$ the exceptional set outside the puncture is exactly the three nonzero two-torsion points. The matrix $-I$ fixes all three, while $R_4$ interchanges $(1/2,0)$ and $(0,1/2)$. Thus the kernel is $\langle-I\rangle$. For $C_6$ the exceptional set consists of the same three two-torsion points and $(1/3,1/3),(2/3,2/3)$. The three two-torsion points force a power of $R_6$ to be $I$ or $-I$, and the latter exchanges the two order-three points. Hence the kernel is trivial. These lists are exhaustive: each nonidentity power has its fixed points given by $(R_n^j-I)v\in\Z^2$; their union is the displayed finite set. The program enumerates these solutions, the kernels, and the matrix centralizers and normalizers. The norm calculation proves exhaustiveness beyond its bounded integer search.
\end{proof}

\begin{claim}{E.24}{Refuted}{Recovery of an arbitrary commuting group from its pointwise subgroup}
For an arbitrary specified $\Gamma\leq C_{\Gamma'}(G)$, its pointwise subgroup is
\[
 \Gamma^{\natural}=\Gamma\cap\Gamma'^{\natural},
\]
and it need not equal $C_{\Gamma'^{\natural}}(G)$. Nor is $\Gamma$ determined as a normalizer by $\Gamma^{\natural}$. For the maximal choice $\Gamma=C_{\Gamma'}(G)$ the proposed equality holds, by the displayed intersection identity.
\end{claim}
\begin{proof}
Both identities follow by testing separately the two defining conditions: membership in $\Gamma$ and pointwise stabilization of each critical datum. To refute recovery, take the first row of E.23. The choices $\Gamma=1$ and $\Gamma=C_2$ have the same pointwise subgroup $1$. Its normalizer inside the centralizer $C_2$ is $C_2$, so the procedure cannot recover the first choice. For the proposed equality without the maximal-choice hypothesis, use the second row and take $\Gamma=1$: its pointwise subgroup is $1$, whereas the centralizer pointwise kernel is $C_2$. All groups and their kernels are included in the finite check. Thus the failure persists even in the specifically requested discrete cases and has a complete intersection formula as its replacement.
\end{proof}
